\section{Conclusion}

We presented \method, a protocol-defined semantic fusion framework for heterogeneous source-coordinate aggregation. Quantum measurement motivates the formalism, but the broader issue is missing-view source fusion: before clients upload prototypes, the protocol must say what makes their statistics comparable. With a public value-and-mask schema, coverage-aware class statistics, and optional high-order sketches, \method\ lets the server fuse evidence across different source coordinates.

Across image experiments, real missing-view benchmarks, quantum-backend compatibility checks, and controlled high-order tasks, correct schema binding and coverage-aware fusion outperform schema-incorrect, zero-fill, shared-view, FedOpt/FedProto, and equal-byte RFF controls under the tested protocols. Strong learned imputation, autoencoder imputation, mask-aware MLP, HeMIS-style fusion, and late fusion define the full-dimensional boundary; they can match or exceed protocol prototypes when trainable reconstruction or missing-modality classifiers are allowed. Under compact communication, schema-aware prototypes remain much stronger than shared-view compression and usually competitive with equal-byte zero-fill fusion. HOP/CHOP is similarly conditional: it is central when evidence is covariance-dominated or cross-source-interaction-driven, while low-order coverage-aware or group-balanced fusion is preferable when first-order evidence is already strong. The contribution is therefore not a universal missing-modality learner, but a protocol argument: semantic coordinate identity, missingness, coverage, and payload should be first-class objects in QFL and heterogeneous-source FL.

\section*{Data and Code Availability}

All experiments are designed to be reproducible from public datasets or generated readout files. MNIST, Fashion-MNIST, CIFAR-10 feature exports, UCI HAR, UCI Multiple Features, MHEALTH, PAMAP2, Hydraulic Systems, WDBC, and Gas Drift boundary-check commands are listed in the supplementary material. The Qiskit Aer and PennyLane studies are circuit-generated readout compatibility checks; no paid IBM hardware data are used. A review source and code archive with fixed seeds, scripts, configuration files, selected result summaries, payload accounting, and cached-feature checksums is submitted as supplementary material. A public repository will be released upon acceptance.

\section*{CRediT Authorship Contribution Statement}

Leyuan Zhang: Conceptualization, Methodology, Software, Validation, Formal analysis, Investigation, Writing -- original draft, Visualization. Yang Long: Supervision, Conceptualization, Methodology, Writing -- review \& editing, Project administration.

\section*{Funding}

Leyuan Zhang acknowledges financial support from the China Scholarship Council (CSC). The funder had no role in the study design, data collection, analysis, interpretation, manuscript preparation, or decision to submit the article for publication.

\section*{Declaration of Competing Interest}

The authors declare that they have no known competing financial interests or personal relationships that could have appeared to influence the work reported in this paper.

\section*{Declaration of Generative AI and AI-Assisted Technologies in the Writing Process}

During the preparation of this work, the authors used AI-assisted tools to support language polishing, editorial organization, and consistency checks. The authors reviewed and edited the content as needed and take full responsibility for the content of the article.
